% ===================================================================
%  કોષ્ટક નં. ૧૩ — હોટેલ. --- ભોજનાલય. કૉફીઘર.
%
%  Traduction, faite en 2026, en gujarati d'aujourd'hui, sur l'ido de
%  texto/io. Regles de la colonne : voir 10-kostak-01.tex. Une seule
%  numerotation. DEUX notes d'auteur, t13-noto-1 en tete et
%  t13-noto-2 posee entre les blocs 8 et 9, comme dans l'ido.
%
%  L'ALINEA 9 EST LE PLUS DIFFICILE DU LIVRET POUR CETTE COLONNE, ET
%  IL FAUT LE DIRE AVANT DE LE LIRE. C'est un menu : crevettes au
%  beurre, tripes a la mode de Caen, poisson, gigot d'agneau, epinards
%  a la creme, fromage, et un Saint-Emilion. Le Gujarat est l'Etat le
%  plus vegetarien de l'Inde — le seul ou la majorite le soit — et le
%  seul grand Etat sous prohibition totale depuis 1960 : l'alcool y
%  est illegal, pas seulement mal vu. La colonne ecrit le menu entier,
%  sans note, sans detour et avec les mots que la langue possede :
%  ઝીંગા les crevettes, આંતરડાં les tripes, માછલી le poisson, જાંઘ le
%  gigot, દારૂ le vin. TROISIEME DEMONSTRATION DE LA MEME LOI : le
%  cochon du tableau 7 gujarati, la charcuterie du tableau 15 haoussa,
%  et ce menu-ci. LA LANGUE A LES MOTS ; C'EST LE FAIRE QUI EST
%  INTERDIT, PAS LE NOMMER. Trois interdits, trois religions, une
%  seule regle.
%
%  ET « ઝીંગા » PORTE LE PARADOXE TOUT ENTIER. La crevette est le mot
%  courant, et la cote gujaratie est l'une des premieres de l'Inde
%  pour la peche debarquee et exportee — vers le Japon, vers l'Europe
%  — pendant qu'elle est l'une des dernieres pour le poisson mange.
%  On peche ici ce qu'on ne mange pas ailleurs qu'ailleurs. Le mot
%  existe donc a plein, dans les bordereaux du port autant que dans la
%  bouche des Kharva de la cote, et il n'y a rien a inventer.
%
%  SEPTIEME REFUS DE LA COLONNE : LE FROMAGE. « પનીર » est le fromage
%  de l'Inde, il est sur toutes les cartes, tout le monde en mange.
%  Mais c'est un caille frais, non sale, non affine, pris au citron et
%  presse le jour meme : ce n'est pas ce qui arrive a la fin d'un
%  repas francais. Le mot serait tombe sous la main comme la
%  chataigne, le coquelicot et la pervenche avant lui. On ecrit
%  « ચીઝ ». QUATRIEME FOIS DE SUITE.
%
%  MAIS « શાક » NE CHANGE PAS DE MOT, IL CHANGE DE PLACE. Le legume
%  d'accompagnement du menu francais porte ici le nom de ce qui est
%  au centre de l'assiette gujaratie — le શાક est le plat, pas la
%  garniture. Le releve n'y touche pas : on ecrit શાક, et c'est le
%  lecteur qui verra que dans ce repas-la il vient apres la viande.
%
%  « શતરંજ » (78) EST LA DEUXIEME EXPORTATION DE SUITE. Au tableau 12
%  c'etait કુલી, le porteur, parti d'ici vers l'anglais ; ici c'est
%  les echecs, chaturanga en sanskrit — les quatre corps d'une armee
%  —, passes par la Perse et par l'arabe avant d'arriver dans le cafe
%  de Royan. Sur la meme table il y a le trictrac et les dames, qu'il
%  a fallu transcrire. Deux tableaux, deux mots qui remontent le
%  courant.
%
%  ENFIN « ધોબણ » REVIENT, ET CETTE FOIS C'EST LA FEMME. Au tableau 9
%  le mot nommait la bergeronnette, l'oiseau qui bat de la queue au
%  bord de l'eau. Ici (42) c'est la blanchisseuse elle-meme, celle
%  pour qui l'oiseau avait ete nomme. Le cercle se ferme a l'interieur
%  de la meme colonne, a quatre tableaux d'intervalle.
% ===================================================================

%%K t13-noto-1 noto
\VUnotes{143.2mm}{%
\textit{\VUgras{(*) ઓરડાની વિગત માટે કોષ્ટક નં. ૬ જુઓ.}}%
}

%%K t13-apar-1 apar
\VUsaut{3.0mm}
\VUtitre{15.0pt}{60}{કોષ્ટક નં. ૧૩}
\VUsaut{1.6mm}
\VUpk{10.2pt}{40}{હોટેલ. --- ભોજનાલય.}
\VUsaut{2.0mm}
\VUpk{10.2pt}{40}{કૉફીઘર.}
\VUsaut{3.4mm}

%%K t13-01-1 p
1. --- ગઈ કાલે સાંજે રોયાં પહોંચીને હું « \textit{મહાસાગર હોટેલ} » માં
ઊતર્યો, જેની ભલામણ મને એક મિત્રે કરી હતી.

%%K t13-01-2 p
મને એક સુંદર \VUgras{ઓરડો} \textsuperscript{(2)} અપાયો, બીજા
\VUgras{માળે} \textsuperscript{(1)}, જ્યાં હું ખૂબ સારી રીતે સૂતો (*).

%%K t13-02-1 p
2. --- આજે સવારે, મારા ઓરડામાંથી નીકળીને, જ્યાં \VUgras{નોકરાણી}
\textsuperscript{(3)} સવારનો નાસ્તો લાવી હતી, હું \VUgras{ધૂમ્રપાનના
ખંડમાં} \textsuperscript{(5)} ગયો, જે \VUgras{વાચનખંડ}
\textsuperscript{(4)} તરીકે પણ વપરાય છે, \VUgras{છાપાં}
\textsuperscript{(6)} વાંચવા માટે, \VUgras{સિગારેટ}
\textsuperscript{(8)} પીતાં પીતાં. વાંચ્યા પછી હું \VUgras{લિફ્ટથી}
\textsuperscript{(9)} નીચે ઊતર્યો અને શહેરમાં લટાર મારવા ગયો.

%%K t13-03-1 p
3. --- હું હોટેલના \VUgras{કારકુનને} \textsuperscript{(12)} પૂછવાનું
ભૂલી ગયો હતો કે \VUgras{સહિયારા ભોજનનું} \textsuperscript{(11)} જમણ કયા
વાગ્યે અપાય છે, એટલે હું વહેલો પાછો આવ્યો અને એનો લાભ લઈને હોટેલના
\VUgras{સ્નાનગૃહમાં} \textsuperscript{(10)} નહાવા ગયો. પછી હું ફરી
\VUgras{ગલ્લે} \textsuperscript{(15)} ગયો, મારા એક મિત્રને ફોન કરવા.
પણ \VUgras{ટેલિફોન} \textsuperscript{(13)} રોકાયેલો હતો ; મારી મૂંઝવણ
જોઈને \VUgras{કૅશિયરે} \textsuperscript{(14)}, જે \VUgras{બિલ}
\textsuperscript{(17)} \VUgras{મુસાફરોના રજિસ્ટર} \textsuperscript{(16)}
પર નોંધતી હતી, \VUgras{દરવાનને} \textsuperscript{(18)} વિનંતી કરી કે
\VUgras{બેલબૉયને} \textsuperscript{(19)} મોકલીને મારું કામ
\VUgras{સાઇકલ} \textsuperscript{(20)} પર કરાવે.

%%K t13-04-1 p
4. --- મેં તેનો આભાર માન્યો અને \VUgras{ભારવાહકને}
\textsuperscript{(24)} (મજૂરને) બક્ષિસ આપવા બહાર નીકળ્યો, જે સ્ટેશનથી
પોતાની \VUgras{હાથલારીમાં} \textsuperscript{(25)} મારી \VUgras{ટોપીની
પેટી} \textsuperscript{(26)}, \VUgras{મારી બૅગ} \textsuperscript{(27)}
અને \VUgras{મારી પ્રવાસની થેલી} \textsuperscript{(28)} લાવ્યો હતો.
હમણાં જ આવેલા \VUgras{ટપાલીએ} \textsuperscript{(21)} મને એક
\VUgras{પત્ર} \textsuperscript{(22)} આપ્યો.

%%K t13-05-1 p
5. --- હું હજી થોડી વાર બારણાના ઉંબરે ઊભો રહ્યો, \VUgras{ટપાલપેટી}
\textsuperscript{(23)} પાસે, રસ્તા પરની અવરજવર જોવા.
\VUgras{કંડક્ટર} \textsuperscript{(30)}, જે \VUgras{ઓમ્નિબસનો}
\textsuperscript{(29)} છે, પોતાની ગાડી પર એક \VUgras{ટ્રંક}
\textsuperscript{(31)} ચડાવતો હતો, જે એક \VUgras{મુસાફરની}
\textsuperscript{(32)} હતી, જેને \VUgras{હોટેલમાલિક}
\textsuperscript{(33)} વળાવવા આવ્યો હતો. એક જુવાન \VUgras{તારવાળો}
\textsuperscript{(35)} નિરાંતે હોટેલના એક ગ્રાહક માટે \VUgras{તાર}
\textsuperscript{(34)} લાવ્યો ; એક \VUgras{પ્રવાસી}
\textsuperscript{(36)}, ખભે \VUgras{કૅમેરો} \textsuperscript{(38)}
લટકાવીને, પોતાની \VUgras{માર્ગદર્શિકા} \textsuperscript{(37)} જોતો હતો.

%%K t13-tit-1 sub
\VUsaut{4.0mm}
\VUpk{10.2pt}{40}{જમવાનો સમય.}
\VUsaut{3.0mm}

%%K t13-06-1 p
6. --- જાળી આગળ એક બેફિકર \VUgras{હમાલ} \textsuperscript{(39)} પોતાના
\VUgras{ઊંચકવાના આંકડા} \textsuperscript{(40)} અને પોતાની
\VUgras{પૉલિશની પેટી} \textsuperscript{(41)} વચ્ચે ઝોકાં ખાતો હતો,
જ્યારે એક જુવાન \VUgras{ધોબણ} \textsuperscript{(42)}, જે કપડાંનો
\VUgras{ટોપલો} \textsuperscript{(43)} ઊંચકતી હતી, \VUgras{ભોજનખંડના
ઉપરીના} \textsuperscript{(44)} ગંભીર ચહેરા પર હસતી હતી, જે બારણા પાસે
ઊભો હતો.

%%K t13-07-1 p
7. --- \VUgras{રસોઇયાને} \textsuperscript{(45)} જોઈને, જે પોતાના
\VUgras{રસોડામાંથી} \textsuperscript{(47)} બહાર આવ્યો હતો — જે
\VUgras{ભોંયરાના માળે} \textsuperscript{(48)} છે — એક \VUgras{રસોડાના
છોકરાને} \textsuperscript{(46)} કામે મોકલવા, મને યાદ આવ્યું કે જમવાનો
સમય પાસે છે, અને હું ભોજનાલયમાં પેઠો, એ ચોકમાં થઈને જ્યાં એક
\VUgras{ડ્રાઇવરે} \textsuperscript{(50)} પોતાની \VUgras{મોટર}
\textsuperscript{(49)} ઊભી રાખી હતી, જ્યારે એક \VUgras{મિકૅનિક}
\textsuperscript{(52)} \VUgras{સાઇકલસવારની} \textsuperscript{(53)}
સૂચના પ્રમાણે \VUgras{પેટ્રોલની ત્રિચક્રી} \textsuperscript{(51)}
સમારતો હતો, જેને હમણાં જ અકસ્માત થયો હતો.

%%K t13-08-1 p
8. --- મેં એક જુવાન \VUgras{બેલબૉયને} \textsuperscript{(54)} પૂછ્યું,
જે \VUgras{બિયરના પ્યાલા} \textsuperscript{(55)} ઉપાડતો હતો, કે સહિયારું
ભોજન વરંડા નીચે છે કે નહીં ; અને તેના હકારમાં જવાબ મળતાં હું એક ટેબલ
પર બેઠો અને મેં મને ખૂબ સરસ જમણ પીરસાવ્યું.

%%K t13-noto-2 noto
\VUnotes{143.2mm}{%
\textit{\VUgras{(*) શબ્દકોશમાં આપેલી વાનગીઓની યાદીની મદદથી શિક્ષક નીચેના
મેનુમાં સહેલાઈથી ફેરફાર કરી શકશે.}}%
}

%%K t13-tit-2 sub
\VUsaut{4.0mm}
\VUpk{10.2pt}{40}{સરસ જમણ.}
\VUsaut{3.0mm}

%%K t13-09-1 p
9. --- વેઇટર મને જમણનું મેનુ (*) અને દારૂની યાદી
લાવ્યો. મેં પસંદ કરીને \VUgras{શરૂઆતની વાનગી તરીકે ઝીંગા}
\VUgras{માખણ} સાથે મગાવ્યાં, અને પહેલી વાનગી તરીકે \VUgras{કૅનની ઢબે
રાંધેલાં આંતરડાં}. મેં \VUgras{માછલી} ન ખાધી, પણ ઘેટાના બચ્ચાની
\VUgras{જાંઘનો} એક ભાગ માગ્યો, અને \VUgras{શાક તરીકે} મલાઈ સાથે
\VUgras{પાલક}. પછી, \VUgras{વચલી વાનગીઓમાંથી} કોઈ મને ન ગમી, એટલે મેં
જુદી જુદી મીઠી વાનગીઓ મગાવી — \VUgras{ચીઝ}, \VUgras{ફળ} અને
\VUgras{બિસ્કિટ}. વળી મેં ઉત્તમ સેંત-એમિલિયોંનો \VUgras{દારૂ}
ઉમેર્યો, અને મારા જમણથી ખૂબ સંતોષ પામીને \VUgras{વેઇટરને}
\textsuperscript{(57)} સારી \VUgras{બક્ષિસ} \textsuperscript{(59)}
આપીને હું ટેબલ પરથી ઊઠ્યો.

%%K t13-10-1 p
10. --- મને નીકળવાની તૈયારીમાં જોઈને \VUgras{વ્યવસ્થાપકે}
\textsuperscript{(60)} મને ઝરૂખે જવાનું સૂચવ્યું, \VUgras{મહાસાગરનું}
\textsuperscript{(62)} ભવ્ય દૃશ્ય જોવા. દૂર \VUgras{માછીમારીની નાની
હોડીઓ} \textsuperscript{(63)}, \VUgras{સઢવાળાં વહાણ}
\textsuperscript{(64)} અને એક પ્રચંડ \VUgras{મુસાફર જહાજ}
\textsuperscript{(65)} દેખાવા લાગ્યાં, જેની કાળી આકૃતિ સફેદ
\VUgras{વાદળ} \textsuperscript{(67)} પર ઊપસી આવતી હતી, જે
\VUgras{ક્ષિતિજનાં} \textsuperscript{(66)} હતાં. હળવી લહેરે
\VUgras{ધજા} \textsuperscript{(70)} ફરકાવી, જે \VUgras{પાટિયાની}
\textsuperscript{(69)} ઉપર ખોડેલી છે, અને એ પાટિયું \VUgras{ખંડનું}
\textsuperscript{(68)} છે.

%%K t13-tit-3 sub
\VUsaut{3.0mm}
\VUpk{10.2pt}{40}{મનોરંજન.}
\VUsaut{3.0mm}

%%K t13-11-1 p
11. --- દૃશ્ય એટલું રસપ્રદ હતું કે મેં નક્કી કર્યું કે કાલે હું
કૉફીઘરની અગાસી પર \VUgras{બેવડું દૂરબીન} \textsuperscript{(71)} કે
\VUgras{દૂરબીન} \textsuperscript{(72)} લઈને ચડીશ.

%%K t13-12-1 p
12. --- ઝરૂખો છોડીને હું હોટેલના કૉફીઘરમાં ગયો, એક \VUgras{પીણું}
\textsuperscript{(73)} પીવા અને \VUgras{બિલિયર્ડની બાજી}
\textsuperscript{(75)} રમવા. હું જરાય થોભ્યા વિના કૉફીના ખંડમાં થઈને
પસાર થયો, જ્યાં ઘણા ગ્રાહકો \VUgras{શતરંજ} \textsuperscript{(78)},
\VUgras{ડ્રાફ્ટ્સ} \textsuperscript{(79)}, \VUgras{ડોમિનો}
\textsuperscript{(80)}, \VUgras{બૅકગૅમન} \textsuperscript{(81)} કે
\VUgras{પત્તાં} \textsuperscript{(82)} રમતા હતા, અને હું સીધો
\VUgras{બિલિયર્ડના ખંડમાં} \textsuperscript{(74)} ચડી ગયો.

%%K t13-13-1 p
13. --- બાજી પૂરી થઈ ત્યારે હું નીચેના માળે ઊતર્યો અને વેઇટર પાસેથી
\VUgras{કવર} \textsuperscript{(84)}, \VUgras{કાગળ} \textsuperscript{(83)},
\VUgras{ટપાલટિકિટ} \textsuperscript{(85)}, \VUgras{કલમ}
\textsuperscript{(87)} અને \VUgras{શાહી} \textsuperscript{(86)} માગ્યાં,
પત્ર લખવા.

%%K t13-13-2 p
પછી મેં એક \VUgras{ગાડીવાનને} \textsuperscript{(92)} બોલાવ્યો, જેની
\VUgras{ગાડી} \textsuperscript{(88)} કૉફીઘર આગળ ઊભી હતી. મેં તેને
કલાકના કે એક ફેરાના ભાવ વિશે પૂછ્યું, અને ભાવ નક્કી થતાં તેણે મારે
માટે \VUgras{ગાડીનું બારણું} \textsuperscript{(89)} ખોલ્યું અને મને
બેસાડ્યો. તે પાછો \VUgras{બેઠક} \textsuperscript{(91)} પર ચડ્યો,
\VUgras{ચાબુકથી} \textsuperscript{(93)} ઘોડાને ઝડપથી ઉપાડ્યા, અને અમે
એક રમણીય લટાર માટે ઊપડ્યા.

\VUsaut{6.0mm}
\VUfilet{22mm}
